\chapter{Introdución}
\label{chap:introducion}

\lettrine{N}ESTE primeiro capítulo da memoria, se explica á motivación, se expoñen os obxectivos e se define á estrutura.

\section{Motivación}
Durante a décadas dos 60 e 70 a industria musical sufriu unha revolución, nunha gran parte dos lugares de restauración comezaron aparecer aparellos que permitían elixir as seguintes cancións no local, as gramolas. Coa dixitalización da música, as gramolas foron desaparecendo e nalgúns casos adaptándose os novos formatos e dixitalizandose. Nos nosos días, coas grandes facilidades que existen e a música na nube é moi complicado atopar un bar no que se permita elixir a seguinte canción, coma antes.

\section{Obxetivo}
O obxetivo principal de esto proxeto é crear unha gramola virtual baseada nun servidor no que tanto o administrador da gramola coma os usuarios dispoñan dunha interface o mais sinxela posible, facendo deste xento mais doada a súa utilización. 

\section{Estado da arte}
\label{Estado da arte}
Existen diferentes tipos de gramolas físicas, similares as antigas pero cun repertorio musical moi amplo e dotadas das maiores novidades. En canto a gramolas sen aparellos físicos controladas polos clientes dende os seus dispositivos, existe un mercado pequeno.

\subsection*{TouchTunes}
Expandida por Norte América e parte de Europa é a mais coñecida, dispón dun gran catálogo, sistema de adaptación da música ao local e diferentes ferramentas para o entretemento dos clientes. Ten un sistema de créditos que se poden consumir nos diferentes establecementos nos que está implantada a plataforma.
\href{https://www.touchtunes.com}{Páxina oficial TouchTunes}

\subsection*{AMI}
Como directa competidora de TouchTunes, AMI entertainment comezou creando pianos nos que introducindo unha moeda podíase elixir algunha das cancións, posteriormente adicaronse a facer gramolas físicas e nos nosos días seguen coa dixitalización destas. Tamén funciona cun sistema de créditos pra facer soar a música que mais che gusta.
\href{https://amientertainment.com}{Páxina oficial de AMI entertainment}

\subsection*{Outloud}
A diferenza das anteriores, esta aplicación está mais pensada pra uso particular, pódese descargar en dispositivos Android ou iOS e a través da rede local os usuarios poden xestionar a lista de reprodución.
\href{https://outloud.dj/}{Páxina oficial Outloud}

\subsection*{Raveberry}
Este proxecto similar a Outloud está pensado pra ser usado dende unha RaspberryPi, o seu punto forte é a integración de multitude de servizos de transmisión de música e permitir os seus usuarios xestionaren a cola de reprodución dun xeito democrácito.
\href{https://github.com/raveberry/raveberry/}{Repositorio de Raveberry en Gitub}

\section{Obxetivos}
O principal obxetivo do proxecto é crear unha aplicación web para xestionar unha gramola virtual. Preténdese que a interación dos usuarios seña o mais sinxela posible, permita crear unha cola de reproducción e escoitala en calqueira dispositivo con \Gls{Spotify}, para isto desenvolverase un \gls{frontend} e un \gls{backend} en \Gls{ionic} e \Gls{django} respectivamente.

% \section{Estrutura da memoria}
% Este documento estruturase en 6 capítulos.
% \begin{itemize}
% 	\item Neste primeiro capítulo.
% 	\item No capítulo 2.
% \end{itemize}
